\section{Enchainer des additions et des soustractions}
\begin{rappelbox}
    Dans un calcul comportant uniquement des additions et des soustractions,
    on peut \bb{transformer toutes les soustractions en additions}.
\end{rappelbox}

\begin{exemplebox}
    \begin{align*}
        & A = 7-(+4)+(-3)-(-5) \\
        & A =\answer{7+(-4)+(-3) + (+5) } \\
        & A =\answer{3+(-3)+(+5)} \\
        & A =\answer{+5}
    \end{align*}
\end{exemplebox}

\begin{methodebox}
\begin{enumerate}
    \item Je transforme les soustractions en additions.
    \item Je supprime les parenthèses lorsque c'est possible.
    \item Je regroupe les nombres positifs et négatifs si cela simplifie
    le calcul.
\end{enumerate}
\end{methodebox}

\begin{exemplebox}

\bb{Etape 1} : Je transforme toutes les soustractions en additions :
\begin{align*}
    A &= -8+3-(-2)+(-5) \\
    A &= \answer{-8+3+(+2)+(-5)} \\
\end{align*}

\begin{minipage}[t]{0.48\textwidth}
    \bb{Méthode 1:} J'applique les règles de priorités opératoires
    \begin{align*}
        A &= \answer{\bb{-8+3}+(+2)+(-5)} \\
          &= \answer{\bb{-5+2}-5} \\
          &= \answer{\bb{-3-5}} \\
          &= \answer{-8}
    \end{align*}
\end{minipage}
\hfill
\begin{minipage}[t]{0.48\textwidth}
    \bb{Méthode 2:} Je regroupe les termes de même signes
    \begin{align*}
        A &= \answer{(\bb{-8-5}) + (\bb{3+2})} \\
          &= \answer{\bb{-13+5}} \\
          &= \answer{-8}
    \end{align*}
\end{minipage}


\end{exemplebox}
